% =========================================================
% mechlib/annotations/dimensions.tex
% BEMASSUNGEN: LINEARES MASS UND WINKELMASS
% =========================================================
% Lineare Bemassungen werden geometrisch aus zwei Endpunkten abgeleitet. Der
% Offset wirkt senkrecht auf die Verbindung von "from" nach "to"; damit
% funktioniert dasselbe Pic fuer horizontale, vertikale und schraege Masse.
%
% Oeffentliche Pics:
%   mech/dimension       lineare Bemassung mit Doppelpfeil
%   mech/angle-dimension Kreisbogen zwischen zwei Winkeln
%
% Interne Koordinaten -A und -B sind die senkrecht versetzten Endpunkte der
% Masslinie. Die Vektorrechnung erfolgt in TikZ "let"-Syntax.
% =========================================================
\pgfkeys{
  /mech/dimension/.is family,
  /mech/dimension,
  from/.initial={(0,0)},
  to/.initial={(1,0)},
  label/.initial={$l$},
  offset/.initial=3mm,
  style/.initial=mech/dimension,
}
\tikzset{
  pics/mech/dimension/.style args={#1}{
    code={
      \pgfkeys{/mech/dimension/.cd,#1}
      \edef\mechDimFrom{\pgfkeysvalueof{/mech/dimension/from}}
      \edef\mechDimTo{\pgfkeysvalueof{/mech/dimension/to}}
      \path let
        \p1=\mechDimFrom,\p2=\mechDimTo,
        \n1={veclen(\x2-\x1,\y2-\y1)},
        \n2={-(\y2-\y1)/\n1*\pgfkeysvalueof{/mech/dimension/offset}},
        \n3={(\x2-\x1)/\n1*\pgfkeysvalueof{/mech/dimension/offset}}
      in
        coordinate (-A) at (\x1+\n2,\y1+\n3)
        coordinate (-B) at (\x2+\n2,\y2+\n3);
      \draw[\pgfkeysvalueof{/mech/dimension/style}]
        (-A)--(-B)
        node[midway,fill=white,inner sep=1pt]
        {\pgfkeysvalueof{/mech/dimension/label}};
    }
  },
  pics/mech/dimension/.default={},
}
\pgfkeys{
  /mech/angle-dimension/.is family,
  /mech/angle-dimension,
  start angle/.initial=0,
  end angle/.initial=45,
  radius/.initial=7mm,
  label/.initial={$\alpha$},
}
\tikzset{
  pics/mech/angle-dimension/.style args={#1}{
    code={
      \pgfkeys{/mech/angle-dimension/.cd,#1}
      \draw
        (\pgfkeysvalueof{/mech/angle-dimension/start angle}:\pgfkeysvalueof{/mech/angle-dimension/radius})
        arc[start angle=\pgfkeysvalueof{/mech/angle-dimension/start angle},
            end angle=\pgfkeysvalueof{/mech/angle-dimension/end angle},
            radius=\pgfkeysvalueof{/mech/angle-dimension/radius}];
      \pgfmathsetmacro{\mechAngleMid}{.5*(\pgfkeysvalueof{/mech/angle-dimension/start angle}+\pgfkeysvalueof{/mech/angle-dimension/end angle})}
      \node at (\mechAngleMid:{.7*\pgfkeysvalueof{/mech/angle-dimension/radius}})
        {\pgfkeysvalueof{/mech/angle-dimension/label}};
    }
  },
  pics/mech/angle-dimension/.default={},
}
